
\def\PUDcodigo{EME108}

\if\COMPILAR0
    \def\PUDcodacad{04507.22}
\else
    \def\PUDcodacad{04506.20}
\fi

\if\COMPILAR0
    \def\PUDdisciplina{Materiais}
\else
    \def\PUDdisciplina{Materiais I}
\fi

\def\PUDsemestre{S4}

\def\PUDhoras{80}

%NOVOS ---------------------------------------------------
\def\PUDhorasTeoria{60}
\def\PUDhorasPratica{20}

\if\COMPILAR0
   \def\PUDinterECA{ \hyperref[PUD:ACE005]{Química Geral (04507.4)} }
\else
   \def\PUDinterECA{ \hyperref[PUD:ACE005]{Química Geral (04506.4)} }
\fi

\def\PUDatuaisECA{---}

\def\PUDrecursos{Projetor multimídia; Quadro branco e pincel; Laboratórios de materiais equipamento com cortadeira metalográfica, embutidora, politrizes, capela, microscópio óptico, máquina universal de ensaios, durômetro e microdurômetro.}

%----------------------------------------------------------

\def\PUDcreditos{4}

\def\PUDprerequisito{ \hyperref[PUD:ACE005]{ACE005} }

\if\COMPILAR0
   \def\PUDpreacad{ \hyperref[PUD:ACE005]{Química Geral (04507.4)} }
\else
   \def\PUDpreacad{ \hyperref[PUD:ACE005]{Química Geral (04506.4)} }
\fi

\def\PUDobjetivos{Compreender a importância dos materiais nas construções de engenharia; Conhecer as propriedades dos materiais e os procedimentos teóricos e práticos de determinação e quantificação das mesmas; Conhecer os meios microestruturais de modificação das propriedades dos materiais.}

\def\PUDementa{Materiais e Engenharia; Estrutura Atômica e Ligações Interatômicas; Estrutura dos Sólidos Cristalinos; Imperfeições nos Sólidos; Difusão; Propriedades Mecânicas dos Metais; Propriedades Elétricas; Propriedades Magnéticas.}

\def\PUDconteudo{ & Ciência e engenharia de materiais. Classifiacação dos materiais. Materais avançados. Necessidade dos materiais modernos. Correlações processamento/estrutura/propriedades/desempenho.
\\
& Estrutura atômica. Ligações atômicas nos sólidos. Forças e energias de ligação. Tipos de ligações atômicas.
\\
& Estruturas cristalinas. Estruturas cristalinas dos metais. Sistemas cristalinos. Pontos, direções e planos cristalográficos. Materiais cristalinos e não-cristalinos.
\\
& Defeitos pontuais. Dicordâncias. Defeitos interfaciais. Defeitos volumétricos. Análises microscópicas.
\\
& Mecanismos de difusão. Difusão em regime estacionário e nãoestacionário. Fatores que influenciam na difusão. Outros caminhos de difusão.
\\
& Conceitos de tensão e deformação. Deformação elástica. Deformação plástica. Propriedades em tração. Tensões e deformações de engenharia. Tensões e deformações verdadeiras. Dureza. Fatores de projeto e segurança.
\\
& Condução elétrica. Lei de Ohm. Condutividade elétrica. Condução eletrônica e iônica. Condução em termos de bandas. Semicondutores. Condução elétrica em cerâmicas iônicas e em polímeros. Comportamento dielétrico. Outras características elétricas dos materiais.
\\
& Conceitos básicos de magnetismo. Diamagnetismo, paramagnetismo e ferromagnetismo. Influência da temperatura sobre o comportamento magnético. Domínios magnéticos e histereses magnéticas. Materiais magnéticos. Supercondutividade.
}

\def\PUDmetodo{Aulas expositivas em que serão abordados conteúdos teóricos através da projeção de slides, desenvolvimentos no quadro e resolução de exercícios práticos e teóricos; Procedimentos práticos de preparação metalográfica (embutimento a quente. Lixamento, polimento e ataque químico), microscopia óptica, ensaios de tração e ensaios de dureza.
}

\def\PUDavalia{As 60 horas da parte teórica da disciplina serão avaliadas na forma de 03 (três) avaliações escritas e teóricas e/ou na forma de trabalhos aplicados, ambos sobre os conteúdos das unidades 1 a 8; As 20 horas da parte prática da disciplina serão avaliadas na forma de construção de documento acadêmico sobre a caracterização de metais a partir de resultados coletados no laboratório de materiais. O documento terá equivalência à 01 (uma) avaliação.}

\def\PUDbibbasica{& \href{www.biblioteca.ifce.edu.br/index.asp?codigo_sophia=24421}{Callister Jr.}, W. D.; Rethwisch, D. G.  Ciência e Engenharia de Materiais - Uma Introdução.  8ª ed.  Rio de Janeiro, RJ: LTC, 2012.  817p.  ISBN: 9788521621249. 
\\ 
& \href{www.biblioteca.ifce.edu.br/index.asp?codigo_sophia=42347}{Shackelford}, J. F.  Ciência dos Materiais.  6ª ed.  São Paulo, SP: Pearson, 2012.  556p.  ISBN: 9788576051602.
\\
& \href{www.biblioteca.ifce.edu.br/index.asp?codigo_sophia=24424}{Silva}, André Luiz V. da Costa e.  Aços e Ligas Especiais.  3ª ed.  São Paulo, SP: Blucher, 2010.  646p.  ISBN: 9788521205180. }
%& CALLISTER JR. , W.  D.;  RETHWISCH, D.  G.;  Ciência e Engenharia de Materiais - Uma Introdução.  8ª edição.  Rio de Janeiro.  Editora LTC.  2012.  817p.  ISBN 9788521621249. 
%\\ 
 %& SILVA, A.  L.  C.;  MEI, P.  R.;  Aços e Ligas Especiais.  3ª edição.  São Paulo.  Editora BLUCHER.  2010.  646p.  ISBN 9788521205180. }

\def\PUDbibcomplementar{ & \href{biblioteca.ifce.edu.br/index.asp?codigo_sophia=23805}{Souza}, S.  A.;  Ensaios Mecânicos de Materiais Metálicos - Fundamentos Teóricos e Práticos.  5ª edição.  São Paulo.  Editora BLUCHER.  1982.  286p.  ISBN 9788521200123.
\\
& \href{www.biblioteca.ifce.edu.br/index.asp?codigo_sophia=44596}{Garcia}, A.;  Spim, J. A.;  Santos, C.  A.;  Ensaios dos Materiais.  2ª ed.  Rio de Janeiro, RJ: LTC, 2012.  365p.  ISBN: 9788521620679.
\\
& \href{biblioteca.ifce.edu.br/index.asp?codigo_sophia=24118}{Colpaert}, Hubertus. Metalografia dos produtos siderúrgicos comuns. 4. ed. São Paulo: Edgard Blücher, 2008. 652 p. + Inclui CD-ROM. Inclui bibliografia e índice. ISBN 9788521204497.
\\
 %& \href{www.biblioteca.ifce.edu.br/index.asp?codigo_sophia=64982}{Ashby}, M.  Seleção de Materiais no Projeto Mecânico.  1ª ed.  Rio de Janeiro, RJ: Elsevier, 2012.  673p.  ISBN 9788535245219. 
 %\\
 %& \href{www.biblioteca.ifce.edu.br/index.asp?codigo_sophia=44594}{Gemelli}, Enori.  Corrosão de materiais metálicos e sua caracterização.  Rio de Janeiro, RJ: LTC, 2001.  183p.  ISBN: 8521612907.
%\\
 & \href{www.biblioteca.ifce.edu.br/index.asp?codigo_sophia=23980}{Chiaverini}, Vicente.  Tecnologia mecânica: estrutura e propriedades das ligas metálicas. vol.1.  2º ed.  São Paulo, SP: Pearson Education do Brasil, 1986.  266p.  ISBN: 0074500899.
\\
& \href{www.biblioteca.ifce.edu.br/index.asp?codigo_sophia=23701}{Chiaverini}, V.  Tecnologia Mecânica: Materiais de Construção Mecânica.  v. 3.  2ª ed.  São Paulo, SP: Pearson Education do Brasil, 1986.  388p.  ISBN: 9780074500910. 
%\\
 %& \href{www.biblioteca.ifce.edu.br/index.asp?codigo_sophia=24061}{Chiaverini}, Vicente.  Tecnologia mecânica: processos de fabricação e tratamento.  vol.2.  2º ed.  São Paulo, SP: Pearson Makron Books, 1986.  315p.  ISBN: 9780074500903.
%\\
 %& \href{www.biblioteca.ifce.edu.br/index.asp?codigo_sophia=59389}{Serra}, Eduardo Torres.  Análise de falhas em materiais utilizados no setor elétrico.  E-book.  Interciência.  554p.  ISBN: 9788571933712.
%\\
 %& \href{www.biblioteca.ifce.edu.br/index.asp?codigo_sophia=58235}{Pavanati}, Henrique Cezar.  Ciência e Tecnologia dos Materiais.  E-book.  Pearson.  196p.  ISBN: 9788543009797. 
 }

\def\PUDautor{Francisco Nélio Costa Freitas}

\def\PUDdata{2013-05-20}

\def\PUDrevisao{1}

\def\PUDrevisor{Samuel Vieira Dias}

\def\PUDdatarevisao{2019-05-15}

\PUDcorpo
